% SPDX-License-Identifier: CC-BY-NC-ND-4.0
% Copyright 2026 Ingolf Lohmann.
% Deterministically generated from ARTICLE_DE.md with Pandoc, then metadata-bound.
% Options for packages loaded elsewhere
\PassOptionsToPackage{unicode}{hyperref}
\PassOptionsToPackage{hyphens}{url}
\PassOptionsToPackage{dvipsnames,svgnames,x11names}{xcolor}
%
\documentclass[
  10pt,
  a4paper,
]{article}
\usepackage{amsmath,amssymb}
\usepackage{iftex}
\ifPDFTeX
  \usepackage[T1]{fontenc}
  \usepackage[utf8]{inputenc}
  \usepackage{textcomp} % provide euro and other symbols
\else % if luatex or xetex
  \usepackage{unicode-math} % this also loads fontspec
  \defaultfontfeatures{Scale=MatchLowercase}
  \defaultfontfeatures[\rmfamily]{Ligatures=TeX,Scale=1}
\fi
\usepackage{lmodern}
\ifPDFTeX\else
  % xetex/luatex font selection
  \setmainfont[]{Latin Modern Roman}
  \setsansfont[]{Latin Modern Sans}
  \setmonofont[]{Latin Modern Mono}
\fi
% Use upquote if available, for straight quotes in verbatim environments
\IfFileExists{upquote.sty}{\usepackage{upquote}}{}
\IfFileExists{microtype.sty}{% use microtype if available
  \usepackage[]{microtype}
  \UseMicrotypeSet[protrusion]{basicmath} % disable protrusion for tt fonts
}{}
\makeatletter
\@ifundefined{KOMAClassName}{% if non-KOMA class
  \IfFileExists{parskip.sty}{%
    \usepackage{parskip}
  }{% else
    \setlength{\parindent}{0pt}
    \setlength{\parskip}{6pt plus 2pt minus 1pt}}
}{% if KOMA class
  \KOMAoptions{parskip=half}}
\makeatother
\usepackage{xcolor}
\usepackage[margin=24mm]{geometry}
\usepackage{color}
\usepackage{fancyvrb}
\newcommand{\VerbBar}{|}
\newcommand{\VERB}{\Verb[commandchars=\\\{\}]}
\DefineVerbatimEnvironment{Highlighting}{Verbatim}{commandchars=\\\{\}}
% Add ',fontsize=\small' for more characters per line
\newenvironment{Shaded}{}{}
\newcommand{\AlertTok}[1]{\textcolor[rgb]{1.00,0.00,0.00}{\textbf{#1}}}
\newcommand{\AnnotationTok}[1]{\textcolor[rgb]{0.38,0.63,0.69}{\textbf{\textit{#1}}}}
\newcommand{\AttributeTok}[1]{\textcolor[rgb]{0.49,0.56,0.16}{#1}}
\newcommand{\BaseNTok}[1]{\textcolor[rgb]{0.25,0.63,0.44}{#1}}
\newcommand{\BuiltInTok}[1]{\textcolor[rgb]{0.00,0.50,0.00}{#1}}
\newcommand{\CharTok}[1]{\textcolor[rgb]{0.25,0.44,0.63}{#1}}
\newcommand{\CommentTok}[1]{\textcolor[rgb]{0.38,0.63,0.69}{\textit{#1}}}
\newcommand{\CommentVarTok}[1]{\textcolor[rgb]{0.38,0.63,0.69}{\textbf{\textit{#1}}}}
\newcommand{\ConstantTok}[1]{\textcolor[rgb]{0.53,0.00,0.00}{#1}}
\newcommand{\ControlFlowTok}[1]{\textcolor[rgb]{0.00,0.44,0.13}{\textbf{#1}}}
\newcommand{\DataTypeTok}[1]{\textcolor[rgb]{0.56,0.13,0.00}{#1}}
\newcommand{\DecValTok}[1]{\textcolor[rgb]{0.25,0.63,0.44}{#1}}
\newcommand{\DocumentationTok}[1]{\textcolor[rgb]{0.73,0.13,0.13}{\textit{#1}}}
\newcommand{\ErrorTok}[1]{\textcolor[rgb]{1.00,0.00,0.00}{\textbf{#1}}}
\newcommand{\ExtensionTok}[1]{#1}
\newcommand{\FloatTok}[1]{\textcolor[rgb]{0.25,0.63,0.44}{#1}}
\newcommand{\FunctionTok}[1]{\textcolor[rgb]{0.02,0.16,0.49}{#1}}
\newcommand{\ImportTok}[1]{\textcolor[rgb]{0.00,0.50,0.00}{\textbf{#1}}}
\newcommand{\InformationTok}[1]{\textcolor[rgb]{0.38,0.63,0.69}{\textbf{\textit{#1}}}}
\newcommand{\KeywordTok}[1]{\textcolor[rgb]{0.00,0.44,0.13}{\textbf{#1}}}
\newcommand{\NormalTok}[1]{#1}
\newcommand{\OperatorTok}[1]{\textcolor[rgb]{0.40,0.40,0.40}{#1}}
\newcommand{\OtherTok}[1]{\textcolor[rgb]{0.00,0.44,0.13}{#1}}
\newcommand{\PreprocessorTok}[1]{\textcolor[rgb]{0.74,0.48,0.00}{#1}}
\newcommand{\RegionMarkerTok}[1]{#1}
\newcommand{\SpecialCharTok}[1]{\textcolor[rgb]{0.25,0.44,0.63}{#1}}
\newcommand{\SpecialStringTok}[1]{\textcolor[rgb]{0.73,0.40,0.53}{#1}}
\newcommand{\StringTok}[1]{\textcolor[rgb]{0.25,0.44,0.63}{#1}}
\newcommand{\VariableTok}[1]{\textcolor[rgb]{0.10,0.09,0.49}{#1}}
\newcommand{\VerbatimStringTok}[1]{\textcolor[rgb]{0.25,0.44,0.63}{#1}}
\newcommand{\WarningTok}[1]{\textcolor[rgb]{0.38,0.63,0.69}{\textbf{\textit{#1}}}}
\usepackage{longtable,booktabs,array}
\usepackage{calc} % for calculating minipage widths
% Correct order of tables after \paragraph or \subparagraph
\usepackage{etoolbox}
\makeatletter
\patchcmd\longtable{\par}{\if@noskipsec\mbox{}\fi\par}{}{}
\makeatother
% Allow footnotes in longtable head/foot
\IfFileExists{footnotehyper.sty}{\usepackage{footnotehyper}}{\usepackage{footnote}}
\makesavenoteenv{longtable}
\setlength{\emergencystretch}{3em} % prevent overfull lines
\providecommand{\tightlist}{%
  \setlength{\itemsep}{0pt}\setlength{\parskip}{0pt}}
\setcounter{secnumdepth}{-\maxdimen} % remove section numbering
\ifLuaTeX
\usepackage[bidi=basic]{babel}
\else
\usepackage[bidi=default]{babel}
\fi
\babelprovide[main,import]{ngerman}
\ifPDFTeX
\else
\babelfont{rm}[]{Latin Modern Roman}
\fi
% get rid of language-specific shorthands (see #6817):
\let\LanguageShortHands\languageshorthands
\def\languageshorthands#1{}
\ifLuaTeX
  \usepackage{selnolig}  % disable illegal ligatures
\fi
\IfFileExists{bookmark.sty}{\usepackage{bookmark}}{\usepackage{hyperref}}
\IfFileExists{xurl.sty}{\usepackage{xurl}}{} % add URL line breaks if available
\urlstyle{same}
\hypersetup{
  pdftitle={Survival of the Anschlussfähigsten},
  pdfauthor={Ingolf Lohmann},
  pdfsubject={Operational-formale Anschlussfähigkeit in Computer-, Informations- und soziotechnischen Systemen},
  pdfkeywords={QIK-VRT, Anschlussfähigkeit, Viabilität, Simulation, formale Verifikation, Lean},
  pdflang={de-DE},
  colorlinks=true,
  linkcolor={blue},
  filecolor={Maroon},
  citecolor={Blue},
  urlcolor={blue},
  pdfcreator={LaTeX via pandoc}}

\title{Survival of the Anschlussfähigsten}
\usepackage{etoolbox}
\makeatletter
\providecommand{\subtitle}[1]{% add subtitle to \maketitle
  \apptocmd{\@title}{\par {\large #1 \par}}{}{}
}
\makeatother
\subtitle{Eine operational-formale Übertragung des evolutionären
Fitnessgedankens auf Computer-, Informations- und soziotechnische
Systeme}
\author{Ingolf Lohmann}
\date{31. Juli 2026}

\begin{document}
\sloppy
\maketitle

\renewcommand*\contentsname{Inhalt}
{
\hypersetup{linkcolor=}
\setcounter{tocdepth}{3}
\tableofcontents
}
\textbf{Autor und Urheber der hier vorgelegten
QIK-VRT-Operationalisierung:} Ingolf Lohmann\\
\textbf{Dokumenttyp:} Wissenschaftliches Grundlagendokument /
formalisierbare Theorie\\
\textbf{Fassung:} 1.0 -- Vorveröffentlichungskandidat\\
\textbf{Datum:} 31. Juli 2026\\
\textbf{Textlizenz:} Creative Commons Namensnennung -- Nicht-kommerziell
-- Keine Bearbeitungen 4.0 International (CC BY-NC-ND 4.0)\\
\textbf{Formalisierungsquellen:} Apache License 2.0\\
\textbf{Publikationsstatus:} Nicht begutachteter Kandidat; kein DOI;
noch nicht auf Zenodo veröffentlicht\\
\textbf{Wahrheitsgrenze:} Alle fünf ausgewiesenen Modellsätze FIT-001,
FIT-002, FIT-003, MAT-001 und MAT-002 wurden im H2-Quellstand am exakten
Branch-Head \texttt{37a946b9eefc21ab369ad56b5fbb1e9c436766e1} im
Push-Lauf \texttt{30627411130} mit Lean 4.19.0 kompiliert, dynamisch auf
Axiome geprüft und jeweils mit leerer Axiomenliste als
\texttt{KERNEL\_VERIFIED} gebunden. Der H3-Zielstand wurde am exakten
Branch-Head \texttt{5196495f07c6f696faf6d23f9cfe353532ac042e} im
erfolgreichen Push-Lauf \texttt{30628327497} erneut mit denselben fünf
leeren Axiomenlisten kernelverifiziert. Der vorhandene
\texttt{KERNEL\_RECEIPT.json} bindet beide Exact Heads und schließt die
Status-Transition von \texttt{FORMAL\_PENDING\_KERNEL} zu
\texttt{FORMAL\_PROVED}. Bewiesen sind ausschließlich die ausdrücklich
definierten abstrakten Modelleigenschaften; weder die biologische
Interpretation noch eine empirische Überlebensprognose folgt daraus.
Repository-Promotion, Zenodo-Publikation, DOI-Zuweisung und systemweite
Fertigstellung sind nicht Teil dieses Receipts und bleiben offen
beziehungsweise \texttt{UNCLAIMED}.

\begin{center}\rule{0.5\linewidth}{0.5pt}\end{center}

\hypertarget{zusammenfassung}{%
\subsection{Zusammenfassung}\label{zusammenfassung}}

Die bekannte Formel „survival of the fittest`` wird häufig irreführend
als „Überleben des Stärksten`` verstanden. In der Evolutionsbiologie
bezeichnet Fitness jedoch keinen absoluten Grad körperlicher Stärke,
sondern einen umweltabhängigen reproduktiven Beitrag, der je nach Modell
absolut, relativ oder durch weitere präzise Fitnessmaße quantifiziert
wird. Der Ausdruck wurde von Herbert Spencer geprägt und später von
Charles Darwin als Bezeichnung für natürliche Selektion übernommen.

Dieses Dokument schlägt für das Computerzeitalter die Formulierung
\textbf{„Survival of the Anschlussfähigsten``} vor. Sie ist ausdrücklich
keine wörtliche Übersetzung und keine neue biologische Definition von
Fitness. Sie ist eine operationale Übertragung des Selektionsgedankens
auf technische und soziotechnische Systeme. Anschlussfähigkeit
bezeichnet dabei die Fähigkeit eines Systems, auf veränderte
Anforderungen mit zulässigen Zustandsübergängen zu antworten, kompatible
Relationen herzustellen und zugleich seine festgelegten Kerninvarianten,
Sicherheitsgrenzen und Evidenzketten zu erhalten.

Der Begriff wird durch beschriftete Übergangssysteme, gültige
Reaktionen, endliche Umgebungsspuren, Simulationsrelationen und
gewichtete Viabilitätsmaße formalisiert. Unter expliziten
Voraussetzungen wird gezeigt: Kann System A jede gültige,
invariantenerhaltende Reaktion von System B simulieren, dann ist die
Menge der von B bewältigbaren Umgebungsspuren in derjenigen von A
enthalten. Unter derselben Verteilung der Umgebungsanforderungen besitzt
A daher keine geringere potentielle Viabilität; bei einer echt größeren
Spurenmenge mit positivem Gewicht besitzt A eine strikt größere
potentielle Viabilität.

Der Satz ist mathematisch, bedingt und maschinenprüfbar. Er behauptet
weder, dass möglichst viele Schnittstellen immer vorteilhaft seien, noch
dass Anschlussfähigkeit allein biologischen Fortpflanzungserfolg
erkläre. Kosten, Sicherheitsrisiken, Fehlentscheidungen,
Ressourcenbegrenzungen und die konkrete Umwelt bleiben
selektionsrelevant. In QIK-VRT bildet Anschlussfähigkeit die Verbindung
von Unterscheidbarkeit, typisierter Information, prüfbarer Relation,
zulässiger Wirkung, Provenienz, Gate-Entscheidung und persistierter
Wirkungsbestätigung.

\textbf{Kanonische Interpretationsregel von Ingolf Lohmann:}\\
\textbf{Survival of the fittest = Survival of the Anschlussfähigsten.}

\textbf{Kernaussage:} Im Computerzeitalter überlebt nicht notwendig das
unveränderteste oder stärkste System. Unter vergleichbaren Bedingungen
bleibt dasjenige System eher fortsetzbar, das neue, relevante Relationen
herstellen kann, ohne dabei seine überprüfbare Identität und seine
Wahrheitsgrenzen zu verlieren.

\begin{center}\rule{0.5\linewidth}{0.5pt}\end{center}

\hypertarget{english-abstract}{%
\subsection{English abstract}\label{english-abstract}}

The phrase ``survival of the fittest'' is often misread as ``survival of
the strongest.'' In evolutionary biology, however, fitness denotes
environment-dependent reproductive contribution, quantified by absolute,
relative, or other precise fitness measures depending on the model,
rather than absolute physical strength. The phrase was coined by Herbert
Spencer and later adopted by Charles Darwin as a description of natural
selection.

This paper proposes \textbf{``survival of the most connectable''} as an
operational transfer of the fitness principle to computer, information,
and sociotechnical systems. It is neither a literal translation nor a
redefinition of biological fitness. Here, \emph{connectability} means
the capability of a system to establish admissible relations and respond
to changing requirements through valid state transitions while
preserving declared core invariants, security boundaries, and evidence
chains.

We formalize this idea using labelled transition systems, valid
responses, finite environment traces, simulation relations, and weighted
viability measures. Under explicit assumptions, if system A can simulate
every valid invariant-preserving response of system B, then every
environment trace viable for B is also viable for A. Relative to the
same distribution of environmental demands, A therefore has no lower
potential viability. The inequality is strict if A accepts additional
traces of positive weight.

The result is conditional, mathematical, and suitable for machine
verification. It does not imply that more interfaces are always better,
that capability guarantees correct policy choice, or that connectability
is empirically identical to biological reproductive fitness. Costs,
attack surface, resource limits, decision policies, and the actual
environment remain relevant. Within QIK-VRT, connectability links
distinguishability, typed information, verifiable relations, admissible
effects, provenance, gate decisions, and persisted effect
acknowledgements.

\begin{center}\rule{0.5\linewidth}{0.5pt}\end{center}

\hypertarget{problemstellung}{%
\subsection{1. Problemstellung}\label{problemstellung}}

Technische Systeme werden nicht allein dadurch dauerhaft nutzbar, dass
sie einen einmal festgelegten Zustand unverändert bewahren.
Betriebssysteme, Netzwerkprotokolle, Datenformate, wissenschaftliche
Modelle, autonome Systeme und Organisationen treffen fortlaufend auf
veränderte Umgebungen. Neue Gegenstellen, Versionen, Messwerte,
Bedrohungen und Anforderungen erzeugen Selektionsdruck.

Ein System kann darauf in mindestens drei grundsätzlich verschiedenen
Weisen reagieren:

\begin{enumerate}
\def\labelenumi{\arabic{enumi}.}
\tightlist
\item
  Es kann jede Veränderung abweisen und dadurch seine momentane Form
  bewahren, aber seine Nutzbarkeit verlieren.
\item
  Es kann jede Veränderung ungeprüft aufnehmen und dadurch Identität,
  Sicherheit oder Wahrheitsstatus verlieren.
\item
  Es kann relevante Veränderungen unterscheiden, zulässige Anschlüsse
  herstellen und dabei festgelegte Invarianten erhalten.
\end{enumerate}

Nur die dritte Möglichkeit wird hier als \textbf{qualifizierte
Anschlussfähigkeit} bezeichnet. Anschlussfähigkeit bedeutet daher weder
grenzenlose Offenheit noch bloße Anzahl von Netzwerkverbindungen. Sie
ist die überprüfbare Fähigkeit zur invariantenerhaltenden Fortsetzung
unter veränderten Bedingungen.

Die zugespitzte Formel lautet:

\[
\boxed{\text{Survival of the fittest}
\quad\rightsquigarrow\quad
\text{Survival of the Anschlussfähigsten}.}
\]

Das Gleichheitszeichen der kanonischen Kurzfassung bezeichnet die
festgelegte Interpretations- und Übersetzungsregel für das
Computerzeitalter. In der wissenschaftlichen Langform wird dafür der
Pfeil \(\rightsquigarrow\) verwendet, um eine falsche lexikalische oder
biologische Identitätsbehauptung auszuschließen.

\begin{center}\rule{0.5\linewidth}{0.5pt}\end{center}

\hypertarget{historische-und-biologische-wahrheitsgrenze}{%
\subsection{2. Historische und biologische
Wahrheitsgrenze}\label{historische-und-biologische-wahrheitsgrenze}}

\hypertarget{urheberschaft-des-ausdrucks}{%
\subsubsection{2.1 Urheberschaft des
Ausdrucks}\label{urheberschaft-des-ausdrucks}}

Herbert Spencer prägte den englischen Ausdruck „survival of the
fittest`` 1864 in \emph{The Principles of Biology}. Alfred Russel
Wallace empfahl Darwin den Ausdruck am 2. Juli 1866. Darwin stimmte am
5. Juli 1866 seiner ergänzenden Verwendung zu, wollte „natural
selection`` aber nicht vollständig ersetzen. Er verwendete die Wendung
1868 in \emph{The Variation of Animals and Plants under Domestication}
und nahm sie 1869 in die fünfte Auflage von \emph{On the Origin of
Species} auf. Das Selektionsprinzip ist darwinisch; die konkrete
Wortfolge stammt von Spencer {[}1--5{]}.

\hypertarget{fitness-ist-nicht-blouxdfes-uxfcberleben}{%
\subsubsection{2.2 Fitness ist nicht bloßes
Überleben}\label{fitness-ist-nicht-blouxdfes-uxfcberleben}}

Biologische Fitness quantifiziert in modernen evolutionstheoretischen
Kontexten den reproduktiven Beitrag eines Organismus oder Genotyps in
einem bestimmten Populations-, Umwelt- und Zeitkontext. Je nach Modell
werden absolute, relative, individuelle, inklusive, zeitlich gemittelte
oder geometrische Fitnessmaße unterschieden. Ein Organismus, der lange
lebt, aber keine erblichen Beiträge zur Folgepopulation leistet, kann in
diesem technischen Sinn eine geringe Fitness besitzen. Umgekehrt kann
ein kurzlebiger Organismus hohen Fortpflanzungserfolg haben {[}6--8{]}.

Deshalb gilt:

\[
\text{biologische Fitness}
\neq
\text{Körperstärke}
\neq
\text{individuelle Lebensdauer}.
\]

Fitness ist außerdem relativ zur jeweiligen Umwelt, Population,
Zeitskala und betrachteten Reproduktionseinheit.

\hypertarget{status-der-vorgeschlagenen-uxfcbersetzung}{%
\subsubsection{2.3 Status der vorgeschlagenen
Übersetzung}\label{status-der-vorgeschlagenen-uxfcbersetzung}}

„Anschlussfähigkeit`` ist in diesem Dokument:

\begin{itemize}
\tightlist
\item
  keine Ersatzdefinition biologischer Fitness;
\item
  keine Behauptung, Darwin habe diesen deutschen Ausdruck verwendet oder
  gemeint;
\item
  keine Aussage, Anschlussfähigkeit sei in jeder Umwelt selektiv
  vorteilhaft;
\item
  kein moralisches oder politisches Werturteil;
\item
  sondern eine operationalisierbare Informatik-Interpretation des
  allgemeinen Gedankens relativer Bewährung unter Selektionsbedingungen.
\end{itemize}

Ohne Population, Replikation, Vererbung, Variation und Selektion ist das
folgende Übergangsmodell keine digitale Evolution im biologischen Sinn.
„Evolvierbarkeit`` als Fähigkeit, vererbbare adaptive Variation
hervorzubringen, ist ebenfalls von der hier definierten
Anschlussfähigkeit zu unterscheiden {[}12--14{]}.

Formal werden daher zwei Größen unterschieden:

\[
w_{\mathrm{bio}}(x,E)
=
\text{biologische Fitness von }x\text{ in Umwelt }E,
\]

\[
a_{\mathrm{tech}}(S,E,H)
=
\text{technische Anschlussfähigkeit von System }S
\text{ in }E\text{ bis Horizont }H.
\]

Es wird ausdrücklich \textbf{nicht} behauptet:

\[
w_{\mathrm{bio}} \equiv a_{\mathrm{tech}}.
\]

\begin{center}\rule{0.5\linewidth}{0.5pt}\end{center}

\hypertarget{grundbegriffe}{%
\subsection{3. Grundbegriffe}\label{grundbegriffe}}

\hypertarget{definition-1-technisches-system}{%
\subsubsection{Definition 1: Technisches
System}\label{definition-1-technisches-system}}

Ein technisches System ist ein Tupel

\[
\mathcal S=(X,x_0,U,\longrightarrow,P,Q),
\]

wobei:

\begin{itemize}
\tightlist
\item
  \(X\) eine Menge möglicher Systemzustände ist;
\item
  \(x_0\in X\) der Anfangszustand ist;
\item
  \(U\) eine Menge unterscheidbarer Umgebungsanforderungen oder Eingaben
  ist;
\item
  \(x\xrightarrow{u}x'\) einen möglichen Zustandsübergang bei
  Anforderung \(u\) bezeichnet;
\item
  \(P:X\to\{\bot,\top\}\) die zu erhaltende Kerninvariante bezeichnet;
\item
  \(Q:X\times U\times X\to\{\bot,\top\}\) den Erfüllungs-, Sicherheits-
  oder Wirkungskontrakt bezeichnet.
\end{itemize}

Vorausgesetzt wird \(P(x_0)=\top\). Beim Vergleich zweier Systeme
bezeichnet ein gemeinsames Symbol \(u\in U\) dieselbe Anforderung nur
dann, wenn eine gemeinsame Vertragssemantik oder eine explizite
Label-Verfeinerungsabbildung dies rechtfertigt. Bloße Gleichheit des
Eingabetyps genügt semantisch nicht.

\hypertarget{definition-2-guxfcltige-reaktion}{%
\subsubsection{Definition 2: Gültige
Reaktion}\label{definition-2-guxfcltige-reaktion}}

Ein Übergang \(x\xrightarrow{u}x'\) ist gültig, wenn

\[
x\xrightarrow{u}x'
\land P(x')
\land Q(x,u,x').
\]

Damit reicht es nicht, dass das System \emph{irgendwie} reagiert. Die
Reaktion muss ausführbar, invariantenerhaltend und vertragsgemäß sein.

\hypertarget{definition-3-umgebungsspur}{%
\subsubsection{Definition 3:
Umgebungsspur}\label{definition-3-umgebungsspur}}

Eine endliche Umgebungsspur ist eine Folge

\[
\sigma=(u_1,u_2,\ldots,u_n)\in U^*.
\]

Ausgehend von \(R_0=\{x_0\}\) wird die Menge gültig erreichbarer
Zustände rekursiv definiert durch

\[
R_{k+1}
=
\left\{
x'\in X
\;\middle|\;
\exists x\in R_k:
x\xrightarrow{u_{k+1}}x'
\land P(x')
\land Q(x,u_{k+1},x')
\right\}.
\]

\hypertarget{definition-4-potentielle-viabilituxe4t}{%
\subsubsection{Definition 4: Potentielle
Viabilität}\label{definition-4-potentielle-viabilituxe4t}}

Das System \(\mathcal S\) bewältigt die Spur \(\sigma\), geschrieben

\[
\operatorname{Viable}(\mathcal S,\sigma),
\]

genau dann, wenn nach Verarbeitung der gesamten Spur mindestens ein
gültiger Zustand erreichbar bleibt:

\[
R_{|\sigma|}\neq\varnothing.
\]

Diese Definition ist existential. Sie beschreibt eine vorhandene
Fähigkeit. Ob eine reale Steuerungsstrategie tatsächlich einen solchen
Pfad auswählt, ist eine zusätzliche Frage.

\hypertarget{definition-5-anschluss-sprache}{%
\subsubsection{Definition 5:
Anschluss-Sprache}\label{definition-5-anschluss-sprache}}

Für einen endlichen Horizont \(N\) ist

\[
L_N(\mathcal S)
=
\{\sigma\in U^{\le N}\mid
\operatorname{Viable}(\mathcal S,\sigma)\}
\]

die Menge aller bis Länge \(N\) bewältigbaren Umgebungsspuren.

\hypertarget{definition-6-relative-anschlussfuxe4higkeit}{%
\subsubsection{Definition 6: Relative
Anschlussfähigkeit}\label{definition-6-relative-anschlussfuxe4higkeit}}

System \(\mathcal A\) ist bis zum Horizont \(N\) mindestens so
anschlussfähig wie System \(\mathcal B\), geschrieben

\[
\mathcal A\succeq_N\mathcal B,
\]

wenn

\[
L_N(\mathcal B)\subseteq L_N(\mathcal A).
\]

Anschlussfähigkeit ist damit kein absoluter Zahlenwert, sondern zunächst
eine Vergleichsrelation innerhalb eines festgelegten Umgebungs- und
Wahrheitsscopes.

\begin{center}\rule{0.5\linewidth}{0.5pt}\end{center}

\hypertarget{konstruktive-simulationsrelation}{%
\subsection{4. Konstruktive
Simulationsrelation}\label{konstruktive-simulationsrelation}}

Die bloße Prüfung sämtlicher Spuren kann exponentiell teuer sein. Eine
lokale Simulationsrelation liefert einen konstruktiven hinreichenden
Beleg.

Seien

\[
\mathcal A=(X_A,a_0,U,\to_A,P_A,Q_A)
\]

und

\[
\mathcal B=(X_B,b_0,U,\to_B,P_B,Q_B)
\]

zwei Systeme über derselben Eingabemenge. Eine Relation

\[
R\subseteq X_B\times X_A
\]

heißt \textbf{invariantenerhaltende Anschluss-Simulation}, wenn:

\begin{enumerate}
\def\labelenumi{\arabic{enumi}.}
\tightlist
\item
  \(R(b_0,a_0)\) gilt;
\item
  für alle \(a,b,u,b'\) gilt: Wenn \(R(b,a)\) und
  \(b\xrightarrow{u}_B b'\) eine gültige Reaktion von \(B\) ist, dann
  existiert ein \(a'\) mit einer gültigen Reaktion
  \(a\xrightarrow{u}_A a'\) und \(R(b',a')\).
\end{enumerate}

Die Richtung ist bewusst gewählt: \(A\) kann jede gültige Reaktion von
\(B\) nachbilden. Dies entspricht der in formalen Methoden üblichen
Verwendung von Simulationsbeziehungen zur Verhaltens- und
Trace-Verfeinerung {[}9--11{]}.

\hypertarget{satz-1-spurerhaltung-durch-anschluss-simulation}{%
\subsubsection{Satz 1: Spurerhaltung durch
Anschluss-Simulation}\label{satz-1-spurerhaltung-durch-anschluss-simulation}}

Simuliert \(A\) das System \(B\) in diesem Sinn, formal
\(\operatorname{ViabilitySimulation}(B,A,R)\), dann gilt für jede
endliche Spur \(\sigma\):

\[
\operatorname{Viable}(\mathcal B,\sigma)
\Longrightarrow
\operatorname{Viable}(\mathcal A,\sigma).
\]

\hypertarget{beweisidee}{%
\paragraph{Beweisidee}\label{beweisidee}}

Induktion über die Länge von \(\sigma\).

\begin{itemize}
\tightlist
\item
  Für die leere Spur folgt die Behauptung aus der Relation der
  Anfangszustände.
\item
  Im Induktionsschritt liefert ein gültiger \(B\)-Übergang für die
  nächste Eingabe aufgrund der Simulationsbedingung einen entsprechenden
  gültigen \(A\)-Übergang. Die Folgezustände stehen erneut in Relation.
  Damit lässt sich jeder endliche gültige \(B\)-Pfad schrittweise durch
  einen gültigen \(A\)-Pfad begleiten.
\end{itemize}

Folglich gilt:

\[
L_N(\mathcal B)\subseteq L_N(\mathcal A)
\quad\text{für jedes }N\in\mathbb N.
\]

\hypertarget{satz-2-die-anschlussrelation-ist-eine-pruxe4ordnung}{%
\subsubsection{Satz 2: Die Anschlussrelation ist eine
Präordnung}\label{satz-2-die-anschlussrelation-ist-eine-pruxe4ordnung}}

Für festes \(N\) ist \(\succeq_N\) reflexiv und transitiv.

\hypertarget{beweis}{%
\paragraph{Beweis}\label{beweis}}

Reflexivität folgt aus

\[
L_N(\mathcal S)\subseteq L_N(\mathcal S).
\]

Transitivität folgt aus der Transitivität der Mengeninklusion:

\[
L_N(\mathcal C)\subseteq L_N(\mathcal B)
\land
L_N(\mathcal B)\subseteq L_N(\mathcal A)
\Longrightarrow
L_N(\mathcal C)\subseteq L_N(\mathcal A).
\]

Antisymmetrie gilt im Allgemeinen nur bis zur beobachtbaren
Spuräquivalenz. Zwei intern verschiedene Implementierungen können
dieselbe Anschluss-Sprache besitzen.

\begin{center}\rule{0.5\linewidth}{0.5pt}\end{center}

\hypertarget{gewichtete-anschlussfuxe4higkeit-und-technisches-uxfcberleben}{%
\subsection{5. Gewichtete Anschlussfähigkeit und technisches
Überleben}\label{gewichtete-anschlussfuxe4higkeit-und-technisches-uxfcberleben}}

Sei \(\Omega_N\subseteq U^{\le N}\) eine endliche Menge relevanter
Umgebungsspuren und sei

\[
\omega:\Omega_N\to\mathbb R_{\ge 0}
\]

eine nichtnegative Gewichtung mit

\[
\sum_{\sigma\in\Omega_N}\omega(\sigma)>0.
\]

Definiere den gewichteten Anschlusswert

\[
A_{N,\omega}(\mathcal S)
=
\frac{
\sum_{\sigma\in L_N(\mathcal S)\cap\Omega_N}\omega(\sigma)
}{
\sum_{\sigma\in\Omega_N}\omega(\sigma)
}.
\]

Wenn die normierten Gewichte als Wahrscheinlichkeitsverteilung
interpretiert werden, ist \(A_{N,\omega}(\mathcal S)\) die
Wahrscheinlichkeit, dass eine gezogene Anforderungsspur
\emph{grundsätzlich} durch mindestens einen gültigen Pfad bewältigt
werden kann.

Das ist eine Wahrscheinlichkeit potentieller Bewältigbarkeit, nicht ohne
Weiteres die tatsächliche Überlebenswahrscheinlichkeit eines laufenden
Systems. Für Letztere braucht es zusätzlich einen Scheduler oder
Controller, eine durch ihn induzierte Verteilung über Entscheidungen und
Fehler oder eine robuste Definition, nach der alle zulässigen
Entscheidungen erfolgreich bleiben.

\hypertarget{satz-3-monotonie-des-anschlusswertes}{%
\subsubsection{Satz 3: Monotonie des
Anschlusswertes}\label{satz-3-monotonie-des-anschlusswertes}}

Aus

\[
\mathcal A\succeq_N\mathcal B
\]

folgt

\[
A_{N,\omega}(\mathcal A)
\ge
A_{N,\omega}(\mathcal B).
\]

\hypertarget{beweis-1}{%
\paragraph{Beweis}\label{beweis-1}}

Wegen

\[
L_N(\mathcal B)\cap\Omega_N
\subseteq
L_N(\mathcal A)\cap\Omega_N
\]

und \(\omega(\sigma)\ge0\) ist die gewichtete Summe über die linke Menge
höchstens so groß wie die gewichtete Summe über die rechte Menge. Beide
Brüche besitzen denselben positiven Nenner.

\hypertarget{satz-4-strikter-anschlussvorteil}{%
\subsubsection{Satz 4: Strikter
Anschlussvorteil}\label{satz-4-strikter-anschlussvorteil}}

Gilt zusätzlich

\[
\sum_{\sigma\in
(L_N(\mathcal A)\setminus L_N(\mathcal B))
\cap\Omega_N}
\omega(\sigma)>0,
\]

dann folgt

\[
A_{N,\omega}(\mathcal A)
>
A_{N,\omega}(\mathcal B).
\]

\hypertarget{korollar-bedingte-technische-uxfcberlebensaussage}{%
\subsubsection{Korollar: Bedingte technische
Überlebensaussage}\label{korollar-bedingte-technische-uxfcberlebensaussage}}

Falls in einem festgelegten technischen Selektionsmodell ein System bis
zum Horizont \(N\) genau dann fortbesteht, wenn die eintretende
Umgebungsspur zu seiner Anschluss-Sprache gehört, dann folgt aus

\[
\mathcal A\succeq_N\mathcal B
\]

eine nicht geringere Fortbestehenswahrscheinlichkeit von \(A\).

Das Korollar ist die präzise mathematische Bedeutung von „Survival of
the Anschlussfähigsten``. Es gilt \textbf{unter den angegebenen
Voraussetzungen}, nicht universell und nicht voraussetzungslos.

\begin{center}\rule{0.5\linewidth}{0.5pt}\end{center}

\hypertarget{warum-mehr-muxf6glichkeiten-nicht-automatisch-besser-sind}{%
\subsection{6. Warum mehr Möglichkeiten nicht automatisch besser
sind}\label{warum-mehr-muxf6glichkeiten-nicht-automatisch-besser-sind}}

Die bisherigen Sätze betreffen die Menge \emph{gültiger} Reaktionen.
Eine zusätzliche Schnittstelle, die Kerninvarianten verletzt, falsche
Semantik transportiert oder eine Sicherheitslücke erzeugt, vergrößert
diese Menge nicht notwendig.

Für reale Systeme müssen mindestens folgende Größen getrennt werden:

\begin{itemize}
\tightlist
\item
  Abdeckung relevanter Anforderungen;
\item
  semantische Korrektheit;
\item
  Erhaltung von Identität und Kerninvarianten;
\item
  Ressourcen- und Energieverbrauch;
\item
  Latenz;
\item
  Angriffsfläche und Schadensrisiko;
\item
  Wiederherstellbarkeit;
\item
  Qualität der Entscheidungsstrategie;
\item
  Reproduzierbarkeit und Provenienz.
\end{itemize}

Ein anwendungsspezifischer Nettoindex kann beispielsweise die Form

\[
A^*(\mathcal S)
=
\alpha C
+\beta F
+\gamma R
+\delta P
-\lambda K
-\rho Z
\]

besitzen, wobei \(C\) die Anforderungsabdeckung, \(F\) die semantische
Vertragstreue, \(R\) die Wiederherstellbarkeit, \(P\) die
Provenienzqualität, \(K\) die Kosten und \(Z\) das Risiko bezeichnen.
Die Gewichte sind kontextabhängig und normativ; dieser Index ist deshalb
kein universelles Naturgesetz.

Insbesondere sind folgende Aussagen falsch oder jedenfalls nicht aus den
obigen Sätzen ableitbar:

\[
\text{mehr Schnittstellen} \Longrightarrow \text{mehr Sicherheit},
\]

\[
\text{mehr mögliche Übergänge} \Longrightarrow \text{bessere reale Entscheidung},
\]

\[
\text{mehr Anschlussfähigkeit in Umwelt }E_1
\Longrightarrow
\text{mehr Anschlussfähigkeit in Umwelt }E_2.
\]

Qualifizierte Anschlussfähigkeit ist \textbf{selektiv, typisiert und
prüfbar}.

\begin{center}\rule{0.5\linewidth}{0.5pt}\end{center}

\hypertarget{nichtfestlegung-als-erhaltung-unterscheidbarer-muxf6glichkeiten}{%
\subsection{7. Nichtfestlegung als Erhaltung unterscheidbarer
Möglichkeiten}\label{nichtfestlegung-als-erhaltung-unterscheidbarer-muxf6glichkeiten}}

Nichtfestlegung bedeutet in diesem Rahmen nicht Beliebigkeit. Sie
bedeutet, mehrere noch nicht widerlegte und semantisch unterscheidbare
Fortsetzungen solange verfügbar zu halten, bis Evidenz oder ein
verbindlicher Vertrag eine Auswahl rechtfertigt.

Sei \(H_t\) die Menge der am Zeitpunkt \(t\) mit der vorliegenden
Evidenz vereinbaren Hypothesen oder Fortsetzungen. Eine vorzeitige
Festlegung ersetzt \(H_t\) ohne hinreichenden Grund durch eine echte
Teilmenge \(H'_t\subsetneq H_t\). Dadurch können später erforderliche
gültige Übergänge verloren gehen.

Die methodische Regel lautet daher:

\[
\text{Unterscheidungen bewahren}
\quad\text{bis Evidenz eine Reduktion trägt.}
\]

Auch hier besteht kein unbedingter Vorteil maximaler Offenheit:
Unmögliche, widerlegte oder sicherheitswidrige Alternativen müssen
ausgeschlossen werden. Anschlussfähigkeit entsteht aus der
kontrollierten Balance zwischen Variation und Selektion.

\begin{center}\rule{0.5\linewidth}{0.5pt}\end{center}

\hypertarget{abbildung-auf-qik-vrt}{%
\subsection{8. Abbildung auf QIK-VRT}\label{abbildung-auf-qik-vrt}}

QIK-VRT kann diese Übertragung operational tragen, weil es Anschluss
nicht nur als syntaktische Kompatibilität, sondern als verantwortete
Wirkungskette modelliert. Die folgende Zuordnung ist eine
Systeminterpretation und muss in der jeweiligen Formalisierung exakt
typisiert werden:

\begin{longtable}[]{@{}
  >{\raggedright\arraybackslash}p{(\columnwidth - 2\tabcolsep) * \real{0.5000}}
  >{\raggedright\arraybackslash}p{(\columnwidth - 2\tabcolsep) * \real{0.5000}}@{}}
\toprule\noalign{}
\begin{minipage}[b]{\linewidth}\raggedright
Ebene
\end{minipage} & \begin{minipage}[b]{\linewidth}\raggedright
Funktion für Anschlussfähigkeit
\end{minipage} \\
\midrule\noalign{}
\endhead
\bottomrule\noalign{}
\endlastfoot
Unterscheidung & Zustände, Versionen, Claims und Gegenstellen bleiben
identifizierbar. \\
Information & Nachrichten besitzen Typ, Semantik, Scope und
Provenienz. \\
Messung & Beobachtungen werden mit Messkontext, Unsicherheit und
Zeitpunkt gebunden. \\
Wirkung & Nicht nur Transport, sondern ein beanspruchter Effekt wird
spezifiziert. \\
Relation & Quelle, Empfänger, Voraussetzung, Ergebnis und Evidenz werden
verbunden. \\
Kausalprüfung & Korrelation, Reihenfolge und autorisierte Wirkung werden
nicht verwechselt. \\
Gate & Unzulässige oder unbelegte Übergänge werden fail-closed
blockiert. \\
Persistenz & Entscheidung, Evidenz und Effect Acknowledgement bleiben
auditierbar. \\
\end{longtable}

Für QIK-VRT bedeutet „anschlussfähiger`` daher nicht: mehr Nachrichten
akzeptieren. Es bedeutet:

\[
\text{mehr relevante Kontexte korrekt verarbeiten}
\]

unter gleichzeitiger Erhaltung von

\[
\text{Integrität}
+\text{Semantik}
+\text{Provenienz}
+\text{Sicherheitsgrenzen}
+\text{Wahrheitsstatus}.
\]

Eine QIK-VRT-konforme Implementierung kann die Definitionen dieses
Dokuments durch folgende Artefakte materialisieren:

\begin{enumerate}
\def\labelenumi{\arabic{enumi}.}
\tightlist
\item
  typisierte Temporal- beziehungsweise Effect-Envelopes;
\item
  explizite Zustands- und Versionsübergänge;
\item
  maschinenlesbare Invarianten;
\item
  Gate-Receipts für zulässige und blockierte Anschlüsse;
\item
  vollständige Erfassung auch fehlgeschlagener Versuche;
\item
  Effect Acknowledgements für später beobachtete Wirkungen;
\item
  append-only Evidenz- und Integritätsketten.
\end{enumerate}

\begin{center}\rule{0.5\linewidth}{0.5pt}\end{center}

\hypertarget{anwendungen}{%
\subsection{9. Anwendungen}\label{anwendungen}}

\hypertarget{protokoll--und-api-evolution}{%
\subsubsection{9.1 Protokoll- und
API-Evolution}\label{protokoll--und-api-evolution}}

Versionierte Protokolle überleben technische Umgebungswechsel eher, wenn
sie Fähigkeiten aushandeln, unbekannte Erweiterungen kontrolliert
behandeln, semantische Invarianten wahren und Fehler explizit machen.
Anschlussfähigkeit ist hier messbar als Menge korrekt verarbeitbarer
Peer- und Versionsspuren.

\hypertarget{verteilte-systeme}{%
\subsubsection{9.2 Verteilte Systeme}\label{verteilte-systeme}}

Knoten müssen wechselnde Latenzen, Partitionen, Wiederholungen,
Reihenfolgen und Softwarestände bewältigen. Anschlussfähigkeit umfasst
nicht nur Erreichbarkeit, sondern die Fähigkeit, unter diesen
Bedingungen konsistente und nachweisbare Wirkungen zu erzeugen.

\hypertarget{cybersicherheit}{%
\subsubsection{9.3 Cybersicherheit}\label{cybersicherheit}}

Ungeprüfte Offenheit erhöht die Angriffsfläche. Eine sicher
anschlussfähige Architektur verbindet daher Erweiterbarkeit mit
Authentisierung, Autorisierung, minimalen Rechten, Typprüfung,
Fail-Closed-Gates und vollständiger Evidenz.

\hypertarget{wissenschaftliche-forschung}{%
\subsubsection{9.4 Wissenschaftliche
Forschung}\label{wissenschaftliche-forschung}}

Modelle, Messgeräte und Datenbestände bleiben wissenschaftlich
anschlussfähig, wenn Einheiten, Unsicherheiten, Kalibrierungen,
Methoden, Versionen und Provenienz erhalten sind. Dadurch können neue
Beobachtungen ältere Ergebnisse prüfen oder neu klassifizieren, ohne
deren Originaldaten umzuschreiben.

\hypertarget{kuxfcnstliche-intelligenz}{%
\subsubsection{9.5 Künstliche
Intelligenz}\label{kuxfcnstliche-intelligenz}}

Ein KI-System ist nicht schon deshalb anschlussfähig, weil es beliebige
Eingaben sprachlich beantwortet. Wissenschaftliche Anschlussfähigkeit
verlangt, dass es Evidenzstatus, Unsicherheit, Quellenbindung,
Gegenmodelle und Aktionsgrenzen bewahrt.

\hypertarget{digital-twins-und-autonome-systeme}{%
\subsubsection{9.6 Digital Twins und autonome
Systeme}\label{digital-twins-und-autonome-systeme}}

Ein Digital Twin muss reale Zustandsänderungen, veraltete Sensorwerte,
Modellabweichungen und sichere Rückwirkungen unterscheiden. Sein
technisches Fortbestehen hängt von der Fähigkeit ab, neue Evidenz
aufzunehmen, ohne Modellzustand und reale Welt unzulässig
gleichzusetzen.

\hypertarget{quantenklassische-laufzeiten}{%
\subsubsection{9.7 Quantenklassische
Laufzeiten}\label{quantenklassische-laufzeiten}}

QPU-Backends unterscheiden sich in Topologie, Kalibrierung, Noise,
Gate-Sätzen und Ergebnisformaten. Eine anschlussfähige Runtime bindet
Backend-, Circuit-, Shot-, Kalibrierungs-, Unsicherheits- und
Effect-Receipts in eine gemeinsame prüfbare Hülle. Sie macht
Quantenresultate nicht deterministisch; sie macht deren Behandlung
nachvollziehbar und wiederholbar.

\hypertarget{organisationen-und-soziotechnische-institutionen}{%
\subsubsection{9.8 Organisationen und soziotechnische
Institutionen}\label{organisationen-und-soziotechnische-institutionen}}

Institutionen bleiben handlungsfähig, wenn sie neue Evidenz und neue
Akteure integrieren können, ohne Zuständigkeit, Verantwortlichkeit und
historische Aufzeichnungen aufzulösen. Diese Analogie ist empirisch zu
prüfen und darf nicht mit einem biologischen Gesetz gleichgesetzt
werden.

\begin{center}\rule{0.5\linewidth}{0.5pt}\end{center}

\hypertarget{falsifizierbarkeit-und-grenzen}{%
\subsection{10. Falsifizierbarkeit und
Grenzen}\label{falsifizierbarkeit-und-grenzen}}

\hypertarget{mathematische-ebene}{%
\subsubsection{10.1 Mathematische Ebene}\label{mathematische-ebene}}

Die formalen Sätze können auf drei Arten scheitern:

\begin{enumerate}
\def\labelenumi{\arabic{enumi}.}
\tightlist
\item
  Eine Definition ist inkonsistent oder unzureichend typisiert.
\item
  Ein Gegenmodell verletzt eine behauptete Folgerung trotz erfüllter
  Voraussetzungen.
\item
  Der angegebene Beweis lässt sich durch den vertrauenswürdigen Kernel
  nicht prüfen.
\end{enumerate}

Ein erfolgreicher Maschinenbeweis zeigt ausschließlich, dass die
Konklusion aus den formalisierten Voraussetzungen folgt.

\hypertarget{empirische-technische-ebene}{%
\subsubsection{10.2 Empirische technische
Ebene}\label{empirische-technische-ebene}}

Die Übertragung kann experimentell geprüft werden. Dazu werden vorab
festgelegt:

\begin{itemize}
\tightlist
\item
  die zu vergleichenden Systeme;
\item
  dieselbe Menge oder Verteilung von Umgebungsspuren;
\item
  derselbe Beobachtungshorizont;
\item
  die zu erhaltenden Invarianten;
\item
  Kosten- und Ressourcenbudgets;
\item
  Sicherheits- und Wirkungskriterien;
\item
  eine Auswertungsregel ohne nachträgliche Auswahl nur erfolgreicher
  Fälle.
\end{itemize}

Die Hypothese lautet dann beispielsweise:

\begin{quote}
Systeme mit nachgewiesen größerer invariantenerhaltender
Anschluss-Sprache besitzen unter derselben Umgebungsspurverteilung und
vergleichbaren Ressourcen eine nicht geringere operative
Fortbestehensrate.
\end{quote}

Eine systematische Verletzung dieser Prognose trotz erfüllter
Voraussetzungen widerlegt das gewählte technische Selektionsmodell oder
zeigt, dass relevante Variablen fehlen.

\hypertarget{biologische-ebene}{%
\subsubsection{10.3 Biologische Ebene}\label{biologische-ebene}}

Dieses Dokument prüft nicht, ob technische Anschlussfähigkeit eine
vollständige biologische Fitnessmetrik ist. Eine solche Behauptung wird
nicht erhoben. Biologische Anwendungen benötigten eigenständige
populationsgenetische Modelle und empirische Daten.

\hypertarget{normative-grenze}{%
\subsubsection{10.4 Normative Grenze}\label{normative-grenze}}

Aus höherer biologischer oder technischer Fitness folgt kein höherer
moralischer Wert. Deskriptive Bewährung und normative Rechtfertigung
sind logisch verschieden.

\begin{center}\rule{0.5\linewidth}{0.5pt}\end{center}

\hypertarget{maschinenpruxfcfbarer-kern}{%
\subsection{11. Maschinenprüfbarer
Kern}\label{maschinenpruxfcfbarer-kern}}

Der formal bestätigte Fünf-Claim-Kern ist in drei voneinander getrennten
Lean-Modulen formuliert. FIT-001, FIT-002, FIT-003, MAT-001 und MAT-002
wurden zunächst im H2-Lauf \texttt{30627411130} am exakten Quell-Head
\texttt{37a946b9eefc21ab369ad56b5fbb1e9c436766e1} und anschließend im
H3-Lauf \texttt{30628327497} am exakten Ziel-Head
\texttt{5196495f07c6f696faf6d23f9cfe353532ac042e} mit Lean 4.19.0
erfolgreich kernelgeprüft. Die maschinenlesbaren Artefakte
\texttt{KERNEL\_EVIDENCE\_H2\_FULL\_PENDING.json} und
\texttt{KERNEL\_EVIDENCE\_H3\_FULL\_TARGET.json} weisen für jede der
fünf Proof-Konstanten eine leere Axiomenliste aus. Der
\texttt{KERNEL\_RECEIPT.json} bestätigt beide Workflows, bindet Quell-
und Ziel-Head und weist die Claim-Transition als geschlossen aus. Alle
fünf Aussagen haben im definierten Modell den Status
\texttt{KERNEL\_VERIFIED}.

\hypertarget{fit-001-endliche-operationale-fortsetzung}{%
\subsubsection{11.1 FIT-001: endliche operationale
Fortsetzung}\label{fit-001-endliche-operationale-fortsetzung}}

\texttt{OperationalContinuation.lean} hält eine biologische
Fitnessstruktur bereits auf Typebene von dem operationalen
Fortsetzungsmodell getrennt. Das technische Modell enthält Zustände,
zulässige Verbindungen und eine lokale Viabilitätsbedingung. Rekursiv
bedeutet \texttt{Survives\ S\ n\ x}, dass vom Zustand \texttt{x} eine
Kette lebensfähiger Anschlüsse der Länge \texttt{n} ausgeht.

Der propositionierte Hauptsatz \texttt{FIT001\_checked} bündelt drei
Aussagen:

\begin{enumerate}
\def\labelenumi{\arabic{enumi}.}
\tightlist
\item
  Fortsetzung über \texttt{n+1} Schritte ist genau ein lebensfähiger
  Anschluss plus Fortsetzung des Nachfolgers über \texttt{n} Schritte.
\item
  Ohne lebensfähigen Nachfolger ist Fortsetzung über einen positiven
  Horizont unmöglich.
\item
  Wer einen längeren endlichen Horizont bewältigt, bewältigt auch jeden
  entsprechend kürzeren Horizont.
\end{enumerate}

\hypertarget{fit-002-anschluss-simulation-und-sprachinklusion}{%
\subsubsection{11.2 FIT-002: Anschluss-Simulation und
Sprachinklusion}\label{fit-002-anschluss-simulation-und-sprachinklusion}}

\texttt{ConnectabilitySimulation.lean} definiert ein typisiertes
beschriftetes Übergangssystem, gültige Schritte, endliche lebensfähige
Spuren und viabilitätserhaltende Simulationen. Die Lean-Quelle enthält
für den Hauptsatz \texttt{FIT002\_checked} einen kernelgeprüften
Beweisterm. Seine belegte Aussage lautet:

\begin{quote}
Simuliert ein Zielsystem jeden gültigen Schritt eines Quellsystems,
erhält dabei Viabilität und deckt alle lebensfähigen Quellzustände ab,
dann enthält seine lebensfähige Sprache diejenige des Quellsystems.
\end{quote}

Formal:

\[
\operatorname{ViabilitySimulation}(B,A,R)
\Longrightarrow
L(B)\subseteq L(A).
\]

Die zusätzliche punktierte Variante bindet die ausgezeichneten
Initialzustände \(b_0\) und \(a_0\) und beweist genau die in Abschnitt 4
verwendete Spracheninklusion. In der Lean-Repräsentation wird der
Wirkungskontrakt \(Q\) in das zulässige Schrittrelationsprädikat
eingefaltet; gemeinsame Labelsemantik bleibt eine explizite
Modellvoraussetzung.

Außerdem ist die daraus definierte globale Relation „mindestens so
anschlussfähig wie`` reflexiv und transitiv. Sie ist somit eine
Präordnung auf beobachtbarem endlichem Verhalten; unterschiedliche
interne Implementierungen können dieselbe lebensfähige Sprache besitzen.

\texttt{WeightedConnectability.lean} bildet die elementare Monotonie
zusätzlich in einem diskreten Kernelmodell ab. Es verwendet ein explizit
endliches, duplikatfreies und horizontbeschränktes Spurenuniversum,
entscheidbare Akzeptanzprädikate sowie natürliche Gewichte. Der
normalisierte Score wird als akzeptierte Gewichtsmasse zusammen mit
demselben strikt positiven Gesamtgewicht repräsentiert. Für einen
gemeinsamen Nenner ist der Vergleich der Quotienten exakt der Vergleich
der Zähler; Division und Gleitkommaarithmetik sind daher nicht Bestandteil
des Beweisterms.

MAT-001 beweist die Monotonie der so repräsentierten Scores unter
Sprachinklusion. MAT-002 beweist strikte Vergrößerung, wenn die größere
Sprache einen konstruktiv im endlichen Universum lokalisierten Trace
akzeptiert, den die kleinere Sprache nicht akzeptiert, und dessen Gewicht
positiv ist. Natürliche Gewichte decken nach gemeinsamer Skalierung
rationale Gewichte ab. Eine Instanziierung für beliebige reelle Gewichte,
maßtheoretische Grenzwerte und empirische Wahrscheinlichkeitsmodelle ist
nicht Teil dieses diskreten Lean-Scopes.

\hypertarget{proof-konstanten}{%
\subsubsection{11.2.1 Proof-Konstanten}\label{proof-konstanten}}

\begin{Shaded}
\begin{Highlighting}[]
\NormalTok{FIT{-}001  QIKVRT.V2.OperationalContinuation.FIT001\_checked}
\NormalTok{FIT{-}002  QIKVRT.V2.ConnectabilitySimulation.FIT002\_checked}
\NormalTok{FIT{-}003  QIKVRT.V2.ConnectabilitySimulation.FIT003\_checked}
\NormalTok{MAT{-}001  QIKVRT.V2.WeightedConnectability.MAT001\_checked}
\NormalTok{MAT{-}002  QIKVRT.V2.WeightedConnectability.MAT002\_checked}
\end{Highlighting}
\end{Shaded}

\hypertarget{geschlossene-h2-h3-exact-head-transition}{%
\subsubsection{11.3 Geschlossene H2--H3-Exact-Head-Transition}%
\label{geschlossene-h2-h3-exact-head-transition}}

Für FIT-001 bis FIT-003 sowie MAT-001 und MAT-002 wurden ausgeführt und
gebunden:

\begin{enumerate}
\def\labelenumi{\arabic{enumi}.}
\tightlist
\item
  vollständiger Lean-Build auf exakt gebundenem Commit;
\item
  keine Platzhalter wie \texttt{sorry} oder \texttt{admit};
\item
  Audit aller zusätzlichen Axiome;
\item
  Prüfung auf ausgeschlossene Proof-Escapes;
\item
  Bindung jedes Satzes an Quelltext, Proof Object und Toolchain-Version;
\item
  reproduzierbarer Build in sauberer Umgebung;
\item
  maschinenlesbare Claim-Matrix;
\item
  SHA-256-Manifest der veröffentlichten Artefakte;
\item
  Trennung zwischen kernel-bewiesenen, konditionalen, empirischen,
  interpretativen und normativen Aussagen;
\item
  weiterhin getrennte, erst danach autorisierte Archivierung und
  DOI-Publikation.
\end{enumerate}

Die Punkte 1 bis 7 und 9 sind für den vollständigen Fünf-Claim-Scope in
beiden Exact-Head-Läufen erfüllt. Der H2-Quell- und Claim-Head
\texttt{37a946b9eefc21ab369ad56b5fbb1e9c436766e1} wurde im Lauf
\texttt{30627411130} kompiliert und für alle fünf Proof-Konstanten mit
leeren Axiomenlisten gebunden. Der H3-Ziel-Head
\texttt{5196495f07c6f696faf6d23f9cfe353532ac042e} wurde im erfolgreichen
Lauf \texttt{30628327497} auf denselben Fünf-Claim-Scope geprüft und
ebenfalls mit fünf leeren Axiomenlisten gebunden.

Der \texttt{KERNEL\_RECEIPT.json} bestätigt die Exact-Head-Bindung beider
Läufe. Proof-Referenzen und Satzformulierungen blieben unverändert; nur
die deklarierte Claim-Klassifikation und der Proof-Status gingen von
\texttt{FORMAL\_PENDING\_KERNEL} beziehungsweise
\texttt{AWAITING\_EXACT\_HEAD\_KERNEL\_RECEIPT} zu
\texttt{FORMAL\_PROVED} beziehungsweise \texttt{KERNEL\_VERIFIED} über.
Eine weitere Ziel-Head-Bestätigung ist für diese Transition nicht
erforderlich; die H2--H3-Transition ist damit geschlossen.

Punkt 8 wird erst für den eingefrorenen Publikationskandidaten
abgeschlossen; Punkt 10 bleibt offen. Der Kernel-Receipt und die
geschlossene H2--H3-Transition liegen vor, aber ein kandidatengebundenes
Zenodo-Dateiset mit finalem SHA-256-Manifest und exakten Upload-Hashes ist
noch nicht veröffentlicht. Erst eine darauf bezogene hashgebundene
Autorisierung darf den Upload freischalten. Der kanonische Reststatus
lautet daher:
\texttt{REPOSITORY\_PROMOTION\_COMPLETE=false},
\texttt{ZENODO\_PUBLISHED=false}, DOI \textbf{offen} und
\texttt{SYSTEM\_WIDE\_COMPLETION=UNCLAIMED}.

Der Lean-Beweis kann den mathematischen Implikationskern abschließen. Er
kann nicht allein die empirische Angemessenheit der Umgebungsverteilung,
die Vollständigkeit der Systemmodellierung oder eine biologische
Identität beweisen.

\begin{center}\rule{0.5\linewidth}{0.5pt}\end{center}

\hypertarget{claim--und-statusmatrix}{%
\subsection{12. Claim- und Statusmatrix}\label{claim--und-statusmatrix}}

\begin{longtable}[]{@{}
  >{\raggedright\arraybackslash}p{(\columnwidth - 4\tabcolsep) * \real{0.3333}}
  >{\raggedright\arraybackslash}p{(\columnwidth - 4\tabcolsep) * \real{0.3333}}
  >{\raggedright\arraybackslash}p{(\columnwidth - 4\tabcolsep) * \real{0.3333}}@{}}
\toprule\noalign{}
\begin{minipage}[b]{\linewidth}\raggedright
ID
\end{minipage} & \begin{minipage}[b]{\linewidth}\raggedright
Aussage
\end{minipage} & \begin{minipage}[b]{\linewidth}\raggedright
Status im Kandidaten
\end{minipage} \\
\midrule\noalign{}
\endhead
\bottomrule\noalign{}
\endlastfoot
HIS-001 & Der Ausdruck „survival of the fittest`` wurde von Herbert
Spencer geprägt und später von Darwin übernommen. & historisch
quellengebunden \\
BIO-001 & Biologische Fitness ist nicht mit Körperstärke oder bloßer
Lebensdauer identisch. & wissenschaftliche Hintergrundannahme \\
TRN-001 & „Survival of the Anschlussfähigsten`` ist die autorenseitig
festgelegte Informatik-Interpretation, keine neue biologische
Definition. & interpretativ / normativ definiert \\
FIT-001 & Positive endliche Fortsetzung erfordert eine Kette
lebensfähiger Anschlüsse; ohne lebensfähigen Nachfolger keine
Fortsetzung. & \texttt{KERNEL\_VERIFIED} im definierten Modell \\
FIT-002 & Eine alle lebensfähigen Startzustände abdeckende
viabilitätserhaltende Simulation impliziert globale Inklusion der
endlichen Viabilitätssprachen. & \texttt{KERNEL\_VERIFIED} im
definierten Modell \\
FIT-003 & Eine initialzustandsgebundene viabilitätserhaltende Simulation
impliziert Inklusion der punktierten endlichen Viabilitätssprachen. &
\texttt{KERNEL\_VERIFIED} im definierten Modell \\
MAT-001 & Im kodierten endlichen Naturgewicht-Modell impliziert
Sprachinklusion Monotonie des gemeinsamen positiv normalisierten
Anschlusswertes. & \texttt{KERNEL\_VERIFIED} im definierten Modell \\
MAT-002 & Im selben Modell impliziert ein konstruktiv lokalisierter
positiver Gewichtsträger der strikten Differenz einen strikt größeren
Anschlusswert. & \texttt{KERNEL\_VERIFIED} im definierten Modell \\
EMP-001 & Es bleibt eine offene empirische Hypothese, dass größere
invariantenerhaltende Anschlussfähigkeit unter vorregistrierten
vergleichbaren Bedingungen eine nicht geringere operative
Fortbestehensrate erzeugt. & offen / empirisch zu prüfen \\
LIM-001 & Ob ein festgelegtes technisches Anschlussmaß empirisch mit
biologischer Fitness oder Evolvierbarkeit korrespondiert, bleibt offen
und wird hier nicht nachgewiesen. & offen; keine Identität behauptet \\
NOR-001 & Größere technische Anschlussfähigkeit oder biologische Fitness
darf für sich allein nicht als moralische Vorzugswürdigkeit behandelt
werden. & normative Schutzregel / deklariert \\
\end{longtable}

\begin{center}\rule{0.5\linewidth}{0.5pt}\end{center}

\hypertarget{schlussfolgerung}{%
\subsection{13. Schlussfolgerung}\label{schlussfolgerung}}

„Survival of the Anschlussfähigsten`` ist wissenschaftlich haltbar, wenn
der Satz als explizit begrenzte Übertragung formuliert wird.

Biologisch bleibt Fitness ein kontextabhängiges Maß des reproduktiven
Beitrags; je nach Modell wird sie absolut, relativ oder in einer anderen
präzisen Form angegeben. Informatikseitig wird Anschlussfähigkeit als
Menge gültig bewältigbarer Umgebungsspuren modelliert. Kann ein System
jede invariantenerhaltende Reaktion eines anderen Systems simulieren,
dann kann es mindestens dieselben endlichen Anforderungsspuren
bewältigen. Unter derselben Gewichtung besitzt es daher keine geringere
potentielle Viabilität.

Damit wird aus einem Aphorismus eine prüfbare Aussage:

\[
\boxed{
L_N(\mathcal B)\subseteq L_N(\mathcal A)
\Longrightarrow
A_{N,\omega}(\mathcal B)
\le
A_{N,\omega}(\mathcal A).
}
\]

Die Aussage bleibt relativ zu Umwelt, Horizont, Invarianten, Vertrag und
Kosten. Sie verherrlicht weder Stärke noch grenzenlose Anpassung. Ein
anschlussfähiges System erhält gerade diejenigen Unterschiede, Grenzen
und Nachweise, die eine verantwortete Fortsetzung ermöglichen.

Die präzise Computerzeitalter-Fassung lautet deshalb:

\begin{quote}
Nicht das unveränderte oder vermeintlich stärkste System setzt sich
notwendig fort. Unter vergleichbaren Selektionsbedingungen besitzt
dasjenige System die größere potentielle Viabilität, das mehr relevante
Veränderungen durch gültige, invariantenerhaltende und überprüfbare
Anschlüsse bewältigen kann.
\end{quote}

Kurz:

\[
\boxed{\text{Survival of the Anschlussfähigsten}.}
\]

\emph{q.e.d. Ingolf Lohmann}

\begin{center}\rule{0.5\linewidth}{0.5pt}\end{center}

\hypertarget{literatur}{%
\subsection{Literatur}\label{literatur}}

\begingroup
\small
\begin{enumerate}
\def\labelenumi{\arabic{enumi}.}
\tightlist
\item
  Spencer, Herbert: \emph{The Principles of Biology}. Band 1. Williams
  and Norgate, London, 1864, S. 444--445. Historischer Scan:
  \url{https://archive.org/download/cu31924003036864/cu31924003036864.pdf}.
\item
  Wallace, Alfred Russel an Charles Darwin, 2. Juli 1866. Darwin
  Correspondence Project, Letter 5140:
  \url{https://www.darwinproject.ac.uk/letter/?docId=letters/DCP-LETT-5140.xml}.
\item
  Darwin, Charles an Alfred Russel Wallace, 5. Juli 1866. Darwin
  Correspondence Project, Letter 5145:
  \url{https://www.darwinproject.ac.uk/letter/?docId=letters/DCP-LETT-5145.xml}.
\item
  Darwin, Charles: \emph{The Variation of Animals and Plants under
  Domestication}. Band 2. John Murray, London, 1868. Darwin Online
  F878.2:
  \url{https://darwin-online.org.uk/converted/published/1868_Variation_F878/1868_Variation_F878.2.html}.
\item
  Darwin, Charles: \emph{On the Origin of Species by Means of Natural
  Selection}. 5. Auflage. John Murray, London, 1869. Darwin Online F387:
  \url{https://darwin-online.org.uk/content/contentblock?basepage=1\&hitpage=1\&itemID=F387\&viewtype=text}.
\item
  Gregory, T. Ryan: ``Understanding Natural Selection: Essential
  Concepts and Common Misconceptions.'' \emph{Evolution: Education and
  Outreach} 2 (2009), S. 156--175. DOI:
  \url{https://doi.org/10.1007/s12052-009-0128-1}.
\item
  Orr, H. Allen: ``Fitness and its role in evolutionary genetics.''
  \emph{Nature Reviews Genetics} 10 (2009), S. 531--539. DOI:
  \url{https://doi.org/10.1038/nrg2603}.
\item
  Wadgymar, Susana M. et al.: ``Defining Fitness in Evolutionary
  Ecology.'' \emph{International Journal of Plant Sciences} 185(3),
  2024, S. 218--227. DOI:
  \url{https://doi.org/10.1086/729360}.
\item
  Lynch, Nancy; Vaandrager, Frits: ``Forward and Backward Simulations:
  I. Untimed Systems.'' \emph{Information and Computation} 121(2), 1995,
  S. 214--233. DOI: \url{https://doi.org/10.1006/inco.1995.1134}.
\item
  Clarke, Edmund M.; Emerson, E. Allen; Sistla, A. Prasad: ``Automatic
  Verification of Finite-State Concurrent Systems Using Temporal Logic
  Specifications.'' \emph{ACM TOPLAS} 8(2), 1986, S. 244--263. DOI:
  \url{https://doi.org/10.1145/5397.5399}.
\item
  Baier, Christel; Katoen, Joost-Pieter: \emph{Principles of Model
  Checking}. MIT Press, 2008.
  \url{https://mitpress.mit.edu/9780262026499/principles-of-model-checking/}.
\item
  Wagner, Günter P.; Altenberg, Lee: ``Complex Adaptations and the
  Evolution of Evolvability.'' \emph{Evolution} 50(3), 1996, S.
  967--976. DOI:
  \url{https://doi.org/10.1111/j.1558-5646.1996.tb02339.x}.
\item
  Adami, Christoph: ``Digital genetics: unravelling the genetic basis of
  evolution.'' \emph{Nature Reviews Genetics} 7 (2006), S. 109--118.
  DOI: \url{https://doi.org/10.1038/nrg1771}.
\item
  Ofria, Charles; Wilke, Claus O.: ``Avida: A Software Platform for
  Research in Computational Evolutionary Biology.'' \emph{Artificial
  Life} 10(2), 2004, S. 191--229. DOI:
  \url{https://doi.org/10.1162/106454604773563612}.
\end{enumerate}
\endgroup

\begin{center}\rule{0.5\linewidth}{0.5pt}\end{center}

\hypertarget{vorgeschlagene-zitation-nach-erfolgter-publikation}{%
\subsection{Vorgeschlagene Zitation nach erfolgter
Publikation}\label{vorgeschlagene-zitation-nach-erfolgter-publikation}}

\begingroup
\small
Lohmann, Ingolf: \emph{Survival of the Anschlussfähigsten: Eine
operational-formale Übertragung des evolutionären Fitnessgedankens auf
Computer-, Informations- und soziotechnische Systeme}. Version 1.0, Jahr
2026. DOI: \textbf{offen}.
\endgroup

\end{document}
